\uuid{W0Io}
\titre{Résolution d'équation différentielle du premier ordre}
\auteur{bodin ; croizat ; sacre}
\organisation{exo7}
\licence{CC BY-NC-SA 4.0}
\datecreate{2022-03-08}
\niveau{L1}
\module{Analyse}
\chapitre{Equation différentielle}
\sousChapitre{Résolution d'équation différentielle du premier ordre}
\difficulte{3}
\theme{}
\exo7qcmid{hilisit.equations-differentielles.053}

\contenu{
\begin{question}
\qclassification{225.01}
Soit l'équation différentielle $2y' = 3y + 1$.
Quelles sont les affirmations vraies ?
\begin{answers}
  \bad{Il y a au moins une solution dont la limite en $-\infty$ est $0$.}
  \good{La solution vérifiant $y(0)=0$ est $y(x) = \frac 13 (e^{\frac32x} - 1)$.}
  \bad{La solution vérifiant $y(0)=0$ est $y(x) = 0$.}
  \bad{La solution vérifiant $y(0)=0$ est $y(x) = e^{\frac32x} - 1$.}
\end{answers}
\begin{explanations}
Notre équation différentielle peut se réécrire sous la forme $y' = \frac 32 y + \frac 12$. Les solutions d'une équation diffférentielle $y' = ay + b$ sont les fonctions $x \mapsto C e^{ax} - \frac ba$. Donc ici les solutions de l'équation différentielle sont les fonctions $f(x) = Ce^{\frac32x} -\frac13$.
La solution vérifiant $y(0)=0$ est $y_0(x) = \frac13e^{\frac32x} -\frac13 = \frac 13 (e^{\frac 32 x} - 1)$. Enfin, puisque la limite de $e^{\frac 32 x}$ en $-\infty$ est nulle, la limite de toutes les fonctions solutions en $-\infty$ sera $-\frac 13$.
\end{explanations}
\end{question}
}
